% Volume III: Probability, Statistics, Information, and Causality
% Continuity note for Chapter 19.
% Previous dependencies: Chapters 13-15 provide random variables, joint and conditional distributions, expectation, conditioning, and Bayes' theorem; Chapters 16-17 provide sampling variability, likelihood, estimators, tests, and model comparison; Chapter 18 provides information and Fisher-information language useful for experiment design.
% New durable concepts: causal questions versus associational questions, structural causal models, directed acyclic graphs, parents, exogenous variables, observational distributions, interventions, graph surgery, do-notation, confounding, backdoor adjustment, front-door intuition, randomized experiments, potential outcomes, average treatment effects, blocking, power, neural perturbation design, and causal representation learning caveats.
% Forward hooks: Later chapters reuse interventions in reinforcement learning, control, neural perturbation experiments, causal representation learning, predictive coding, and scientific model comparison.
% Notation commitments: Observed variables are uppercase X,Y,Z,M,S,R; realized values are lowercase x,y,z,m,s,r; interventions use P(Y=y\mid \operatorname{do}(X=x)); structural equations use V_i=f_i(\operatorname{Pa}_i,U_i); potential outcomes are Y_i(1),Y_i(0); treatment assignment is A_i; average treatment effect is \tau.
% Figure plan: no figures are included in this authored source. Proposed raster figures are listed in imagegen-figures/ch19-figure-metadata.md.
\section{Chapter 19: Causality, Interventions, and Experimental Design}
\label{ch:causality-interventions-experimental-design}

Probability describes what is uncertain. Statistics estimates unknown quantities from data. Causality asks a different question: what would change if we acted on the system? A neuron may fire more on trials with large pupil diameter, but that association alone does not say whether pupil-linked arousal caused the firing change, whether task difficulty caused both, or whether an unmeasured state drove the whole pattern.

The central problem of this chapter is to separate association from intervention. We need a language for causal assumptions, a notation for actions, and experimental designs that make causal claims defensible. The point is not to replace probability. Causal models add extra structure to probability models: they say how variables are generated and how that generation changes when one mechanism is externally set.

\paragraph{Chapter problem.}
How can we state, identify, estimate, and interpret causal effects in neuroscience and machine learning without confusing observational correlation with the result of an intervention?

\paragraph{Section map.}
Section 19.1 introduces structural causal models, directed graphs, confounding, interventions, and the backdoor adjustment formula. Section 19.2 gives an intuition-first view of do-calculus and works through a front-door-style identification mechanism. Section 19.3 develops randomized experiments, potential outcomes, treatment effects, neural interventions, and design tradeoffs. Section 19.4 previews causal representation learning, emphasizing invariance, interventions, and identifiability limits.

\subsection{19.1 Correlation, causation, and structural models}

\paragraph{Guiding question.}
When does a statistical relationship such as \(P(Y\mid X=x)\) answer the causal question \(P(Y\mid \operatorname{do}(X=x))\), and when does it fail?

\paragraph{Beginner-facing intuition.}
Correlation is about what tends to be observed together. Causation is about what would happen under an action. If high firing rate and successful behavior are observed together, that is an association. If externally increasing the firing rate would improve behavior, that is a causal claim. The two agree only under assumptions about how the variables were produced.

A causal model is a story with equations. It says that each variable is generated from its direct causes plus its own unexplained noise. A directed graph records which variables are direct inputs to which mechanisms. An intervention changes one mechanism: instead of allowing \(X\) to be generated by its usual causes, we set \(X=x\) from outside the system.

\paragraph{Formal definitions and notation.}
A directed acyclic graph, or DAG, is a directed graph with no directed cycles. Its vertices are variables. If \(Z\to X\), then \(Z\) is a parent of \(X\). Write \(\operatorname{Pa}_i\) for the parents of variable \(V_i\).

A structural causal model is a tuple
\[
M=(G,\{f_i\}_{i=1}^d,P_U),
\]
where \(G\) is a DAG on variables \(V_1,\ldots,V_d\), \(U=(U_1,\ldots,U_d)\) are exogenous noise variables with distribution \(P_U\), and each observed variable is generated by a structural equation
\[
V_i=f_i(\operatorname{Pa}_i,U_i).
\]
The induced observational distribution is the joint law \(P(V_1,\ldots,V_d)\) obtained by sampling \(U\) and evaluating the equations.

An intervention \(\operatorname{do}(X=x)\) replaces the structural equation for \(X\) by the constant assignment
\[
X:=x.
\]
Graphically, this deletes arrows pointing into \(X\) while leaving arrows out of \(X\) in place. The interventional distribution \(P(Y\mid \operatorname{do}(X=x))\) is the distribution of \(Y\) in this modified model.

A confounder for the effect of \(X\) on \(Y\) is a variable \(Z\) that causally influences both \(X\) and \(Y\), producing an open backdoor path
\[
X \leftarrow Z \to Y.
\]
Conditioning or adjusting for appropriate confounders can sometimes recover an interventional effect from observational data.

\paragraph{Core mechanism: backdoor adjustment.}
Suppose \(Z\) blocks all noncausal backdoor paths from \(X\) to \(Y\), and \(Z\) is not a descendant of \(X\). Then
\[
\boxed{
P(Y=y\mid \operatorname{do}(X=x))
=\sum_z P(Y=y\mid X=x,Z=z)P(Z=z).
}
\]
For continuous \(Z\), the sum becomes an integral.

Here is the mechanism in the simplest confounded graph \(Z\to X\), \(Z\to Y\), and \(X\to Y\). Under \(\operatorname{do}(X=x)\), the intervention cuts the arrow \(Z\to X\), but it does not change how \(Z\) is distributed. Therefore
\[
P(Y=y\mid \operatorname{do}(X=x))
=\sum_z P(Y=y\mid \operatorname{do}(X=x),Z=z)P(Z=z).
\]
If \(Z\) blocks the backdoor paths, then once \(X=x\) and \(Z=z\) are fixed, the conditional response mechanism for \(Y\) is the same in the observational and interventional worlds:
\[
P(Y=y\mid \operatorname{do}(X=x),Z=z)
=P(Y=y\mid X=x,Z=z).
\]
Substituting gives the adjustment formula. The formula is not a statistical identity; it depends on the causal graph assumptions.

\paragraph{Concrete worked example.}
Let \(Z\in\{0,1\}\) indicate task difficulty, with \(Z=0\) easy and \(Z=1\) hard, and suppose \(P(Z=1)=1/2\). Let \(X=1\) mean a neural stimulation was applied, and let \(Y=1\) mean success on the trial. Suppose stimulation is more often given on hard trials:
\[
P(X=1\mid Z=0)=0.2,\qquad P(X=1\mid Z=1)=0.8.
\]
The conditional success probabilities are
\[
P(Y=1\mid X=1,Z=0)=0.9,\qquad P(Y=1\mid X=1,Z=1)=0.4,
\]
\[
P(Y=1\mid X=0,Z=0)=0.8,\qquad P(Y=1\mid X=0,Z=1)=0.2.
\]
The causal effect compares difficulty-balanced interventions:
\[
P(Y=1\mid \operatorname{do}(X=1))
=0.5(0.9)+0.5(0.4)=0.65,
\]
\[
P(Y=1\mid \operatorname{do}(X=0))
=0.5(0.8)+0.5(0.2)=0.50.
\]
The causal effect is positive: \(0.65-0.50=0.15\).

The observational association can point the other way. Since stimulation is common on hard trials,
\[
P(Z=1\mid X=1)=0.8,\qquad P(Z=1\mid X=0)=0.2.
\]
Thus
\[
P(Y=1\mid X=1)=0.2(0.9)+0.8(0.4)=0.50,
\]
while
\[
P(Y=1\mid X=0)=0.8(0.8)+0.2(0.2)=0.68.
\]
The observational comparison suggests stimulation is harmful, even though the adjusted causal calculation says it helps. The discrepancy is confounding by difficulty.

\paragraph{Computational neuroscience example.}
Suppose a cortical oscillation \(X\) and a memory score \(Y\) are positively correlated. A latent arousal state \(Z\) may increase both oscillatory power and performance. A structural model can represent \(Z\to X\), \(Z\to Y\), and possibly \(X\to Y\). Observing the correlation \(P(Y\mid X)\) does not determine whether stimulating the oscillation would change memory. If arousal is measured well enough and satisfies the backdoor assumptions, adjustment over \(Z\) is one route to an estimate. A direct stimulation experiment is stronger, but even then off-target effects and state dependence must be modeled.

\paragraph{AI and machine-learning example.}
In a recommender system, \(X=1\) may mean an item was shown prominently, and \(Y=1\) may mean the user clicked. User interest \(Z\) affects both whether the system chooses to show the item and whether the user clicks. The observational probability \(P(Y=1\mid X=1)\) overestimates the effect of promotion if high-interest users are preferentially exposed. A causal effect asks what click probability would result if the same population of users were randomly assigned to see or not see the item.

\paragraph{Common misconceptions and failure modes.}
\begin{itemize}[leftmargin=*]
\item A large correlation is not a causal effect. It may be produced by a common cause, selection bias, reverse causation, or measurement artifacts.
\item Conditioning on every available variable is not always safe. Conditioning on a collider, such as \(X\to C\leftarrow Y\), can create artificial association.
\item A DAG is not learned from data alone without assumptions. Observational distributions can often be compatible with multiple causal graphs.
\item The phrase ``\(X\) causes \(Y\)'' is incomplete without specifying the intervention, population, outcome, and time scale.
\item A structural equation is a modeling assumption, not a claim that biology literally computes the equation.
\end{itemize}

\paragraph{Exercises.}
\begin{enumerate}[label=\textbf{19.1.\arabic*.}, leftmargin=*]
\item \textbf{Mechanical.} In the graph \(Z\to X\), \(Z\to Y\), \(X\to Y\), list the parents of each variable and identify the backdoor path from \(X\) to \(Y\).
\item \textbf{Proof or derivation.} Starting from the law of total probability under \(\operatorname{do}(X=x)\), derive the backdoor adjustment formula for a discrete confounder \(Z\), stating the causal assumption used in the second equality.
\item \textbf{Numerical/computational.} Using the probabilities in the worked example, compute \(P(Y=1\mid X=1)-P(Y=1\mid X=0)\) and compare it with the adjusted causal effect.
\item \textbf{Interpretation.} In a neural stimulation study, why might task difficulty confound the relationship between stimulation and behavior if stimulation is applied more often after difficult trials?
\item \textbf{Misconception/debugging.} A student says, ``We controlled for many variables, so the effect must be causal.'' Give two reasons this conclusion may still fail.
\end{enumerate}

\subsection{19.2 Do-calculus intuition}

\paragraph{Guiding question.}
How can graph surgery and conditional independence rules turn expressions with \(\operatorname{do}(\cdot)\) into observational quantities?

\paragraph{Beginner-facing intuition.}
Do-calculus is a rulebook for replacing interventional expressions by observational expressions when the graph justifies the replacement. The rules are not magic algebra. Each rule says that, after cutting or deleting certain arrows, a conditional independence in the modified graph makes one observation or intervention irrelevant.

The practical mental model is this: first draw the causal graph, then modify it according to the interventions, then read off which paths are blocked. If the graph says a variable can no longer carry information to the outcome after the surgery, it can be removed or exchanged with an observation.

\paragraph{Formal definitions and notation.}
For a DAG \(G\), write \(G_{\bar X}\) for the graph obtained by deleting all arrows into \(X\). This represents the intervention \(\operatorname{do}(X=x)\). Write \(G_{\underline X}\) for the graph obtained by deleting all arrows out of \(X\). Conditional independence in a graph is written
\[
Y\perp\!\!\!\perp Z\mid W \quad \text{in }G,
\]
meaning that all relevant paths from \(Y\) to \(Z\) are blocked by \(W\) according to d-separation.

The three do-calculus rules can be stated informally as follows. Let \(X\) denote variables already intervened on, \(Z\) variables being manipulated or observed, and \(W\) variables being conditioned on.
\begin{enumerate}[label=\textbf{Rule \arabic*.}, leftmargin=*]
\item \textbf{Insert or delete observations.} If \(Z\) is conditionally independent of \(Y\) given \(X,W\) in the graph with incoming arrows to \(X\) removed, then observing \(Z\) gives no extra information about \(Y\) after \(X,W\).
\item \textbf{Exchange action and observation.} If, after the appropriate graph surgery, setting \(Z\) and observing \(Z\) have the same consequences for paths into \(Y\), then \(\operatorname{do}(Z=z)\) may be replaced by conditioning on \(Z=z\).
\item \textbf{Insert or delete actions.} If, after surgery, the intervention on \(Z\) cannot affect \(Y\) given \(X,W\), then \(\operatorname{do}(Z=z)\) can be removed.
\end{enumerate}
The formal statements are compact but easy to misuse. In this course the important skill is not memorizing graph subscripts; it is checking the modified graph before moving a term across the \(\operatorname{do}\)-operator.

\paragraph{Core mechanism: front-door identification.}
The backdoor adjustment fails when the treatment \(X\) and outcome \(Y\) share an unobserved confounder \(U\). Sometimes a measured mediator \(M\) still identifies the effect. Consider the graph
\[
U\to X,\qquad U\to Y,\qquad X\to M\to Y,
\]
with no direct arrow \(X\to Y\), no unblocked backdoor path from \(X\) to \(M\), and all backdoor paths from \(M\) to \(Y\) blocked by conditioning on \(X\). Under these front-door conditions,
\[
\boxed{
P(Y=y\mid \operatorname{do}(X=x))
=
\sum_m P(M=m\mid X=x)
\sum_{x'} P(Y=y\mid M=m,X=x')P(X=x').
}
\]

The mechanism has two steps. First, because there is no backdoor path from \(X\) to \(M\), the effect of setting \(X\) on the mediator is visible observationally:
\[
P(M=m\mid \operatorname{do}(X=x))=P(M=m\mid X=x).
\]
Second, the effect of setting \(M=m\) on \(Y\) can be adjusted by averaging over \(X\), because conditioning on \(X\) blocks the confounding path from \(M\) back through \(X\) and \(U\) to \(Y\):
\[
P(Y=y\mid \operatorname{do}(M=m))
=\sum_{x'}P(Y=y\mid M=m,X=x')P(X=x').
\]
Combining the mediator distribution with the mediator effect gives the front-door formula.

\paragraph{Concrete worked example.}
Let \(X\in\{0,1\}\) be a stimulation command, \(M\in\{0,1\}\) an observed circuit activation, and \(Y\in\{0,1\}\) a behavioral response. Suppose
\[
P(X=1)=0.5,
\]
\[
P(M=1\mid X=0)=0.2,\qquad P(M=1\mid X=1)=0.8,
\]
and
\[
P(Y=1\mid M=0,X=0)=0.1,\quad P(Y=1\mid M=0,X=1)=0.3,
\]
\[
P(Y=1\mid M=1,X=0)=0.6,\quad P(Y=1\mid M=1,X=1)=0.8.
\]
For \(m=0\),
\[
\sum_{x'}P(Y=1\mid M=0,X=x')P(X=x')
=0.5(0.1)+0.5(0.3)=0.2.
\]
For \(m=1\),
\[
\sum_{x'}P(Y=1\mid M=1,X=x')P(X=x')
=0.5(0.6)+0.5(0.8)=0.7.
\]
Therefore
\[
P(Y=1\mid \operatorname{do}(X=1))
=0.2(0.2)+0.8(0.7)=0.60,
\]
whereas
\[
P(Y=1\mid \operatorname{do}(X=0))
=0.8(0.2)+0.2(0.7)=0.30.
\]
The mediator calculation identifies a causal effect of \(0.30\) under the front-door assumptions.

\paragraph{Computational neuroscience example.}
Suppose sensory stimulus \(X\) affects an intermediate neural population \(M\), which affects behavior \(Y\), while an unmeasured alertness variable \(U\) influences both stimulus sampling and behavior. If all stimulus effects on behavior pass through the measured population \(M\), and if the remaining assumptions are plausible, a front-door-style analysis can use the measured population as a mediator. The conclusion is only as strong as the assumptions: unmeasured parallel pathways from \(X\) to \(Y\) or hidden confounders of \(M\) and \(Y\) would break the formula.

\paragraph{AI and machine-learning example.}
In an online platform, \(X\) may be a ranking policy, \(M\) a user engagement signal after seeing the ranked page, and \(Y\) a later purchase. Unmeasured user intent may affect both the chosen ranking and the purchase. If engagement mediates the ranking effect and the required graphical conditions are defensible, front-door reasoning suggests how to estimate a policy effect from logs. In practice, platforms often combine such modeling with randomized experiments because the assumptions are hard to verify.

\paragraph{Common misconceptions and failure modes.}
\begin{itemize}[leftmargin=*]
\item Do-calculus does not make causal effects identifiable from any graph. Some graphs leave effects unidentified.
\item The rules are about conditional independences in modified graphs, not ordinary algebraic cancellation.
\item A mediator is not automatically useful. It must satisfy graphical conditions, and measuring it with noise can introduce additional problems.
\item Front-door identification does not say the mediator is the only biologically meaningful mechanism; it says the measured mediator is sufficient under a specific graph.
\item Replacing \(\operatorname{do}(X=x)\) by \(X=x\) without checking a graph is exactly the mistake do-calculus was designed to prevent.
\end{itemize}

\paragraph{Exercises.}
\begin{enumerate}[label=\textbf{19.2.\arabic*.}, leftmargin=*]
\item \textbf{Mechanical.} For the graph \(Z\to X\), \(X\to Y\), \(Z\to Y\), draw \(G_{\bar X}\) in words by listing which arrows remain after \(\operatorname{do}(X=x)\).
\item \textbf{Proof or derivation.} Derive the front-door formula by writing \(P(Y\mid \operatorname{do}(X=x))=\sum_m P(Y\mid \operatorname{do}(X=x),M=m)P(M=m\mid \operatorname{do}(X=x))\) and stating the graphical replacement used for each factor.
\item \textbf{Numerical/computational.} Recompute the worked front-door example if \(P(M=1\mid X=1)=0.7\) while all other numbers stay fixed. What is \(P(Y=1\mid \operatorname{do}(X=1))\)?
\item \textbf{Interpretation.} In a neural circuit study, what biological assumption is represented by the front-door condition that all directed paths from \(X\) to \(Y\) pass through \(M\)?
\item \textbf{Misconception/debugging.} A student says, ``If \(M\) is between \(X\) and \(Y\), we should always condition on \(M\) to estimate the effect of \(X\).'' Explain why conditioning on a mediator answers a different question from the total causal effect.
\end{enumerate}

\subsection{19.3 Randomized experiments and neural interventions}

\paragraph{Guiding question.}
Why does random assignment support causal estimation, and what extra design issues arise when the intervention is a neural perturbation?

\paragraph{Beginner-facing intuition.}
Randomization breaks the link between treatment assignment and the uncontrolled causes of the outcome. If treatment is assigned by a coin flip, then treated and control groups should be similar on average before treatment, including on variables we did not measure. This makes a simple difference in outcomes interpretable as a causal effect, up to sampling variability and design limitations.

In neuroscience, randomization is necessary but not sufficient. A neural intervention may have off-target effects, change behavior through discomfort or attention, interact with brain state, or be measured with delay and noise. Good experimental design pairs random assignment with controls, blocking, calibration, outcome predefinition, and humility about what intervention was actually delivered.

\paragraph{Formal definitions and notation.}
Let units be indexed by \(i=1,\ldots,n\). A unit may be an animal, a trial, a neuron, or a session, depending on the design. Let \(A_i\in\{0,1\}\) be the treatment assignment, where \(A_i=1\) means treated. The potential outcomes are
\[
Y_i(1),\qquad Y_i(0),
\]
where \(Y_i(1)\) is the outcome that would be observed for unit \(i\) under treatment and \(Y_i(0)\) under control. Only one is observed:
\[
Y_i=A_iY_i(1)+(1-A_i)Y_i(0).
\]
The average treatment effect, or ATE, is
\[
\tau=\mathbb E[Y(1)-Y(0)]
=\mathbb E[Y\mid \operatorname{do}(A=1)]-\mathbb E[Y\mid \operatorname{do}(A=0)].
\]
Random assignment means
\[
A_i\perp\!\!\!\perp (Y_i(1),Y_i(0),Z_i),
\]
where \(Z_i\) denotes pre-treatment covariates. In words, assignment is independent of potential outcomes and pre-treatment features.

\paragraph{Core theorem: difference in means is unbiased under random assignment.}
Suppose \(n_1\) units are assigned to treatment and \(n_0\) to control by a randomized design, and assume no interference between units. The difference-in-means estimator is
\[
\widehat{\tau}
=\frac{1}{n_1}\sum_{i:A_i=1}Y_i
-\frac{1}{n_0}\sum_{i:A_i=0}Y_i.
\]
Under random assignment,
\[
\mathbb E[\widehat{\tau}]=\mathbb E[Y(1)-Y(0)]=\tau,
\]
where the expectation is over the random assignment and sampling of units.

The reason is simple. In the treated group, the observed outcome is \(Y_i(1)\), and random assignment makes the treated group an unbiased sample of the population's treatment potential outcomes. Thus
\[
\mathbb E\left[\frac{1}{n_1}\sum_{i:A_i=1}Y_i\right]=\mathbb E[Y(1)].
\]
Similarly, the control mean estimates \(\mathbb E[Y(0)]\). Subtracting gives \(\tau\). Blocking or stratification preserves this logic within strata and can reduce variance when the blocking variable predicts the outcome.

\paragraph{Concrete worked example.}
A study randomly assigns \(40\) trials to stimulation and \(40\) to sham. The mean behavioral score is \(0.62\) in stimulated trials and \(0.54\) in sham trials. The estimated ATE is
\[
\widehat{\tau}=0.62-0.54=0.08.
\]
If the sample variances are \(s_1^2=0.09\) and \(s_0^2=0.07\), an approximate standard error is
\[
\operatorname{SE}(\widehat{\tau})
=\sqrt{\frac{s_1^2}{40}+\frac{s_0^2}{40}}
=\sqrt{0.004}=0.063.
\]
The estimate is positive, but the standard error is large relative to the effect. The experiment suggests a possible effect, not a precise conclusion. More trials, lower noise, better blocking, or a larger effect would be needed for a sharper estimate.

\paragraph{Computational neuroscience example.}
In an optogenetic experiment, \(A_i=1\) might mean a laser pulse targets a specific cell type on trial \(i\), and \(Y_i\) might be choice accuracy or reaction time. A well-designed experiment includes sham stimulation, random trial order, control of light artifacts, pre-specified outcome windows, and possibly blocking by baseline brain state or stimulus condition. The causal claim is not ``this brain area causes behavior'' in the abstract. It is a claim about the effect of this perturbation protocol on this outcome in this population and context.

\paragraph{AI and machine-learning example.}
An A/B test randomly assigns users to model \(A=0\) or model \(A=1\) and compares click-through rate, retention, or calibration. Randomization protects against user-interest confounding that would affect observational logs. Sequential deployment adds complications: interference between users, novelty effects, nonstationary traffic, and adaptive stopping can bias naive analyses if ignored.

\paragraph{Common misconceptions and failure modes.}
\begin{itemize}[leftmargin=*]
\item Randomization removes systematic confounding in expectation, not random imbalance in every finite sample.
\item A statistically significant effect is not automatically scientifically important; effect size and uncertainty matter.
\item A neural perturbation is an intervention on an apparatus and tissue state, not a pure intervention on a single mathematical variable.
\item Post-treatment variables should not be adjusted for when estimating a total effect unless the estimand is explicitly changed.
\item Repeated trials from the same animal or session may not be independent units; hierarchical structure affects standard errors.
\end{itemize}

\paragraph{Exercises.}
\begin{enumerate}[label=\textbf{19.3.\arabic*.}, leftmargin=*]
\item \textbf{Mechanical.} Define \(A_i\), \(Y_i(1)\), \(Y_i(0)\), observed \(Y_i\), and \(\tau\) for a two-condition neural stimulation experiment.
\item \textbf{Proof or derivation.} Show that if \(A\perp\!\!\!\perp (Y(1),Y(0))\), then \(\mathbb E[Y\mid A=1]-\mathbb E[Y\mid A=0]=\mathbb E[Y(1)-Y(0)]\).
\item \textbf{Numerical/computational.} A randomized experiment has treatment mean \(12.4\), control mean \(10.9\), treatment variance \(9\), control variance \(7\), and \(n_1=n_0=50\). Compute the estimated effect and approximate standard error.
\item \textbf{Interpretation.} Why might blocking by stimulus difficulty improve precision in a neural perturbation experiment?
\item \textbf{Misconception/debugging.} A student analyzes only trials where the animal responded quickly, after treatment was applied. Explain why this post-treatment selection can bias a causal estimate.
\end{enumerate}

\subsection{19.4 Causal representation learning preview}

\paragraph{Guiding question.}
What would it mean for a learned representation to capture causal factors rather than merely predictive correlations, and why is this hard?

\paragraph{Beginner-facing intuition.}
A representation is a map from raw observations to features. A useful predictive representation may encode whatever correlations help on the training distribution. A causal representation aims to encode stable factors that continue to behave predictably under interventions, changes of environment, or distribution shifts.

The hard part is that latent causes are usually not directly observed. A video frame may be generated by object identity, pose, lighting, camera angle, and background. Many different latent descriptions can produce the same observed distribution. Without interventions, multiple environments, temporal structure, or additional assumptions, the causal coordinates are often not identifiable.

\paragraph{Formal definitions and notation.}
Let \(X\in\mathcal X\) be an observation, \(Y\in\mathcal Y\) a target, and
\[
\phi:\mathcal X\to\mathcal H
\]
a representation map into a feature space \(\mathcal H\). A latent-variable generative model writes
\[
C=(C_1,\ldots,C_k),\qquad X=g(C,N),
\]
where \(C_j\) are latent causal factors and \(N\) is nuisance noise. A causal representation would ideally recover features \(H=\phi(X)\) whose coordinates align, perhaps up to simple transformations, with causal factors \(C_j\) relevant for interventions and prediction.

Environments are indexed by \(e\in\mathcal E\). The observational law in environment \(e\) is \(P_e(X,Y)\). An invariance principle says that if \(S\) contains the direct causal predictors of \(Y\), then
\[
P_e(Y\mid C_S)
\]
may remain stable across environments even when the marginal distribution \(P_e(C_S)\) or nuisance correlations change. In representation learning, one tries to find \(\phi\) such that
\[
P_e(Y\mid \phi(X))
\]
is stable across training environments and predictive in new environments. This is an assumption-guided objective, not a guarantee from observational data alone.

\paragraph{Core mechanism: invariance versus spurious prediction.}
Consider binary label \(Y\in\{0,1\}\), causal feature \(C=Y\), and spurious feature \(S\). In environment \(e=1\),
\[
P(S=Y)=0.9,
\]
while in environment \(e=2\),
\[
P(S=Y)=0.6.
\]
The conditional law \(P(Y\mid C)\) is stable: \(Y=C\). The conditional law \(P(Y\mid S)\) changes across environments because the correlation between \(S\) and \(Y\) changes. A classifier that uses \(S\) can perform well in the first environment and fail under shift. A classifier that uses \(C\) generalizes across these shifts.

The mathematical lesson is that predictive accuracy on one distribution does not identify causal features. Stability across environments can provide a clue, but only if the environments perturb nuisance correlations in informative ways and leave the target mechanism stable.

\paragraph{Concrete worked example.}
Suppose the training environment has
\[
P(Y=1)=1/2,\qquad P(S=Y)=0.95,
\]
and the test environment has
\[
P(S=Y)=0.50.
\]
A classifier \(\hat Y=S\) has training accuracy \(0.95\), but test accuracy \(0.50\). A classifier \(\hat Y=C\), where \(C\) is the true causal feature, has accuracy \(1\) in both environments. If \(C\) is not observed, learning it from \(X\) requires assumptions about how \(X\) is generated, access to environments that vary \(S\), interventions, temporal information, or supervision.

\paragraph{Computational neuroscience example.}
Neural population activity \(X\) may reflect stimulus \(C_1\), movement \(C_2\), arousal \(C_3\), and recording drift \(N\). A decoder trained in one task condition may rely on movement-related activity even if the scientific question concerns stimulus encoding. Across conditions where movement correlations change, an invariant stimulus representation should remain predictive of stimulus while the movement shortcut fails. This does not prove the brain explicitly represents a causal graph; it is a modeling strategy for separating stable task variables from nuisance covariation.

\paragraph{AI and machine-learning example.}
In image classification, background may correlate with object class in the training data. A model can learn background texture instead of object shape. Causal representation learning tries to use augmentations, counterfactual data, multiple environments, weak labels, or interventions to learn features tied to the object-generating factors rather than dataset-specific shortcuts. In reinforcement learning, representations of controllable state variables are especially valuable because policies act through interventions on the environment.

\paragraph{Common misconceptions and failure modes.}
\begin{itemize}[leftmargin=*]
\item A disentangled-looking representation is not automatically causal. Axes can be visually interpretable without supporting interventions.
\item High out-of-distribution accuracy suggests robustness but does not by itself identify causal factors.
\item Observational data from one environment usually cannot distinguish causal features from stable but noncausal correlations.
\item Invariance assumptions can fail when the target mechanism itself changes across environments.
\item Biological interpretation should be separated from modeling utility: a useful latent factor in an analysis need not be a literal neural variable.
\end{itemize}

\paragraph{Exercises.}
\begin{enumerate}[label=\textbf{19.4.\arabic*.}, leftmargin=*]
\item \textbf{Mechanical.} Define a representation map \(\phi:\mathcal X\to\mathcal H\) and name the domain, codomain, input, and output in an image or neural-data example.
\item \textbf{Proof or derivation.} In the binary spurious-feature setup, derive the accuracy of the classifier \(\hat Y=S\) in an environment where \(P(S=Y)=q\).
\item \textbf{Numerical/computational.} If \(q=0.9\) in training and \(q=0.55\) in testing, compute the train and test accuracy of \(\hat Y=S\). Compare with a classifier using the causal feature \(C=Y\).
\item \textbf{Interpretation.} In a neural decoding study, how could a representation that encodes movement instead of stimulus still achieve high stimulus-decoding accuracy in one dataset?
\item \textbf{Misconception/debugging.} A student says, ``The network learned a causal representation because the latent dimensions are disentangled.'' Explain what additional evidence or assumptions would be needed.
\end{enumerate}

\paragraph{Closing bridge.}
Causality adds action-oriented structure to probability. Structural models say how variables are generated, do-notation says what changes under intervention, and randomized designs make causal comparisons estimable with fewer assumptions. The next volume turns to optimization: once we can state objectives, causal estimands, and experimental losses, we need mathematical tools for minimizing, constraining, and interpreting them.

\clearpage
